\section{Weighted side computations with binary phase frames}
\label{sec:complex-compression}
We retain the complex motif's data banks, central operations, three tensor
stages, and endpoint corrections. Scalars lie in $\mathbb Z[i,1/2]$ and
labels in $F=\mathbb F_2^h$ with its standard dot product. Throughout this
construction $h\ge8$ is even. Write $\mathcal T=\binom{[h]}3$,
$v=\binom h3$, $N=v^3$ and $m=h^3$.

For triples $S,T$, the central operation has coefficient
$(|S\cap T|-1)/2$. Its side correction has coefficient $+1/2$ for
disjoint triples and $-1/2$ for intersection-two triples. Intersection-one
pairs need no side correction; the diagonal central coefficient is one.
We compress these two side relations separately. All coefficients below
are characteristic-zero coefficients; binary labels do not turn the scalar
arithmetic into XOR arithmetic.

\subsection{A reusable weighted circuit transfer}
An all-ones rectangle with $a$ sources and $b$ targets uses $a+b-1$ roles.
Allocate $a$ input slots and $b$ output slots sharing one pivot. Add the
other inputs into the pivot, then add it into the other outputs. This
invertible mixer $L$ has an all-ones output-by-input submatrix. Its inverse
undoes these additions in reverse order. A fixed dyadic rectangle weight
is placed in the output injection $J$, not in the binary labels.

More generally, a circuit of disjoint-support binary sums with $c$ addition
nodes and $q$ output uses has an invertible embedding on $c+q$ roles.
Allocate one outgoing role per use of each input. At an addition, reuse
its first incoming role as one outgoing pivot, allocate the other outgoing
roles freshly, add the second input to the pivot, and fan it out. The
second incoming role retires. There are $2c+q$ uses and $c$ role
identifications, proving the count. Starting from zero side values,
disjoint input supports give the exact integer coefficients, not merely
their reductions modulo two. Equal sums may have shared fanout.

Let $V$ copy data into the input slots, let $J$ inject weighted outputs,
and retain the central gather $G$ and scatter $R$. The signed chronological
schedule is
\[
 L,\ -J,\ L^{-1},\ -R,\ V,\ G,\ R,\ L,\ J,\ L^{-1},\ -G,\ -V.
\]
An arbitrary side vector $z$ contributes $-JLz+JL(z+Vx)=JLVx$.
The two inverse mixers and final negative copy restore $z$. Similarly the
centers contribute $RGx$ and return to their original values. Thus the
invocation is $y\gets y+x$ precisely when $JLV+RG=I$. Reversing and
inverting the signed operations implements $y\gets y-x$.

For the frame transfer use the original decomposition $A=B\perp P$,
where $P$ is the earlier tensor-indicator line, and suppress the future
norm-one tensor line $Q$. Put
\[
 D_0=B\otimes F,\quad D_1=A\otimes F,\quad
 D_U=(B\otimes F)\perp(P\otimes U),
\]
\[
 X_t=D_{\langle t\rangle},\qquad Y_t=D_{t^\perp}.
\]
The forward schedule's common gate frames, in the displayed order, are
\[
 D_0,D_0,D_0,D_0,X_t,D_1,D_0,D_U,Y_t,D_1,D_1,D_1.
\]
Copy and injection gates are grouped by data triple. In a rectangle the
middle mixer has one common $U$. In a sum circuit it has one $U$ per node.
For the inverse invocation with logical banks exchanged, the frames are
\[
 D_0,D_0,D_0,X_t,D_U,D_1,D_0,Y_t,D_1,D_1,D_1,D_1.
\]
The inverse middle mixer uses physical-$X$ target supports reachable from
each forward node. These grow along reversed edges; forward source supports
must not be substituted for them.

A valid assignment must make every side-role transition increasing and
nondegenerate, with each nonzero orthogonal residual nonalternating. Merely
having nested source sets is insufficient over $\mathbb F_2$. The following
two constructions explicitly meet this stronger condition in both directions.

\subsection{Disjoint triples: coordinate enclosures}
Partition the ordered disjoint-triple relation into rectangles
$\mathcal S\times\mathcal U$. Their ground-point unions are disjoint.
In a forward middle mixer choose its frame $U$ as follows:
\begin{itemize}
\item If $\mathcal S=\{S\}$, use $U=\langle t_S\rangle$.
\item Otherwise, if $\mathcal U=\{T\}$, use $U=t_T^\perp$.
\item Otherwise, use the coordinate subspace on $\bigcup\mathcal S$.
\end{itemize}
For the inverse invocation interchange the two families in this rule.
In every case $\langle t_S\rangle\subset U\subset t_T^\perp$ for
every rectangle pair. These spaces are nondegenerate. In the two singleton
cases the only new nontrivial local complement is that of two orthogonal
triple lines. It has a norm-one coordinate vector outside their six-point
union, since $h>6$.

In the third case both unions have at least four points. The complement
of any source triple line within $U$ contains a coordinate unit outside
that triple. The complement of $U$ within $t_T^\perp$ contains a coordinate
unit in the target union outside $T$. The spaces $U$ and $U^\perp$ are
coordinate spaces. Consequently all nonzero residuals in
$0\subset\langle t_S\rangle\subset U\subset t_T^\perp\subset F$,
and all skipped transitions actually used by input-only or output-only
roles, have norm-one vectors. Every such nondegenerate residual has an
orthonormal basis by the retained binary-space lemma.

The partition is explicit. For disjoint $a$- and $b$-subsets of $[n]$,
$0\le a,b\le3$, allow row stars, column stars, and a split into $l$ and
$n-l$ points. Split cases by the source and target cardinalities in the
first part, and use the product of the already chosen smaller partitions.
For a plan let $(C,L,R)$ be its rectangle count and source/target membership
counts. Products multiply these counts coordinatewise and cases add them.
An empty relation has zero counts; if one subset size is zero, use the full
rectangle. Try all split sizes and choose the least $L+R-C$, breaking ties
by row stars, then column stars, then increasing split size. This finite
recursion is a construction, not a global rectangle optimality claim.

At $h=26$, the counts are
\[
 (C,L,R)=(24870,217255,299571),\qquad R_0=L+R-C=491956.
\]
The exact checker verifies that each of the $4604600$ ordered disjoint
triple pairs appears once. It also checks the frame-case hypotheses for
every generated rectangle. Every such rectangle has scalar weight $+1/2$.

\subsection{Intersection two: an orthonormal sum compiler}
Fix the unique common pair $P=\{i,j\}$. The source and target triples
are $P\cup\{a\}$ and $P\cup\{b\}$ with $a\ne b$ outside $P$.
The binary vectors
\[
 t_a=e_i+e_j+e_a\quad(a\notin P)
\]
are pairwise orthonormal. Since $h$ is even, adjoining
\[
 r_i=e_i+\sum_{a\notin P}e_a,\qquad
 r_j=e_j+\sum_{a\notin P}e_a
\]
gives an orthonormal basis of the entire $F$: each extra vector has odd
weight $h-1$, their mutual product is $h-2=0$ in $\mathbb F_2$, and both
are orthogonal to every $t_a$.

Compute each leave-one-out sum with prefix and suffix sums on the $n=h-2$
inputs. Omit unused totals and intern equal input-support sets. The active
circuit has $3n-6$ additions and $n$ outputs, hence $4n-6$ roles. Each sum
is cancellation-free over the integers. Inject each output with weight
$-1/2$. Distinct triples intersecting in two points have a unique common
pair, so these circuits give every required coefficient exactly once.

At a forward node let $U$ span its ancestor vectors $t_a$; in the inverse
middle circuit let $U$ span the reachable output vectors $t_b$. Nested
support differences are spanned by subsets of an orthonormal family.
No source can reach its excluded output. At the final transition into
the relevant $t_b^\perp$, the residual has an orthonormal basis consisting
of the remaining $t_a$ and the two extra vectors $r_i,r_j$. Cleanup
complements have the same property. Thus every required residual is
zero or nonalternating and nondegenerate in both directions.

This argument applies to any disjoint-support circuit for the same relation,
not only prefix/suffix sums. It is a reusable phase-frame transfer criterion.
At $h=26$ the role count is
\[
 R_2=\binom{26}{2}(4\cdot24-6)=29250,
 \qquad R_{\rm side}=R_0+R_2=521206.
\]

\section{Whole-network composition and accounting}
All local frames above are tensored with the original future line $Q$.
The original data boundaries remain
$A\otimes\langle t\rangle\otimes Q$ to $D_1\otimes Q$ on $X$ and
$B\otimes\langle t\rangle\otimes Q$ to $Y_t\otimes Q$ on $Y$.
Their residuals and the original tensor complements have norm-one witnesses
as in the retained proof. Tensoring nonalternating nondegenerate factors
preserves a norm-one witness. All new side transitions are increasing.
Only the $h+1$ centers decrease once per invocation, by dimension $h$.
Thus $L=3v^2h(h+1)$ is unchanged.

Pair consecutive ground points and put $\pi(T)=\{x\mathbin\oplus1:x\in T\}$.
This is an involution of triples. A triple with a full ground-point pair
intersects its image in two points; otherwise the intersection is empty.
Hence $t_T\perp t_{\pi(T)}$. Match the complete stage-1 auxiliary bank
at fixed triples $(A,B)$ to the stage-3 bank at $(B,\pi(A))$, preserving
each local role index. These banks have equal size, the matching is bijective,
and no simultaneous roles are identified. Each invocation restores arbitrary
scratch, so its scalar correctness is preserved.

All auxiliary first frames are $D_0\otimes Q$ and all last frames are
$D_1\otimes Q$. The joining labels are therefore
\[
 E=F\otimes\langle t_A\rangle\otimes\langle t_B\rangle,
 \qquad H=\langle t_B\otimes t_{\pi(A)}\rangle^\perp\otimes F.
\]
They are nondegenerate and $E\subset H$. The residual $H\cap E^\perp$
contains a norm-one coordinate tensor whose second coordinate lies outside
$A\cup\pi(A)$; such a coordinate exists because $h>6$. Thus the join
also has an orthonormal residual basis. Replacing the old edges
$E\to F^{\otimes3}$ and $0\to H$ by $E\to H$ saves exactly $m$
residual dimensions per removed role and introduces no decrease.

Consequently the entire network has
\[
 W=2N+2v^2(R_{\rm side}+h+1),\qquad s=Wm-2N+2L.
\]
Forward, inverse, forward stages send $(X,Y)$ to $(-Y,X)$ and restore every
surviving auxiliary role. Data endpoint labels and all scratch endpoints
$0,F^{\otimes3}$ are unchanged. The retained weight-27 phase corrections,
signed-bank correction and return permutation therefore give the same
$C^{\otimes m}$ operator on every role, for any number of columns. Each
orthonormal residual contributes precisely its dimension in forward or
inverse translation kernels. These are the full complex-interface obligations.

At $h=26$,
\[
 m=17576,\quad v=2600,\quad N=17576000000,\quad L=14236560000,
\]
\[
 W=7082222160000,\quad s=124477130005280000,
 \quad Wm-s=6678880000,\quad \eta=\frac{19}{354111108}.
\]
Exact logarithm bounds give $\eta>(5/10^9)\log m$, and hence
$s/W<m^{1-5/10^9}$. Thus this construction supports
$a_{\rm c}=1-\sigma=5\cdot10^{-9}$. The exact saving is approximately
$5.48945\cdot10^{-9}$. The finite even-$h$ screen is not an all-$h$ ceiling.

\section{Conditional assembly and precision}
Keep the published bit saving $a_{\rm b}=296/10^{11}$ and the supplied
compact-control movement proof. Choose
\[
 \epsilon=199/1000,\quad c=1,\quad\beta=1/100,\quad
 \zeta=1/1000,\quad\delta=1/10000,\quad C_1=4961/1000,
\]
\[
 \lambda=1-293/10^{11},\qquad\lambda'=1-29/10^{10},\qquad
 \kappa=2^{-31}.
\]
Here $\sigma<\tau$, so the internal exponent is $\chi=\tau$.
The leaf saving is $(1-\beta)a_{\rm c}>a_{\rm b}$ and the reservation
exponent is zero. The seven retained margins have exact minimum
\[
 G_*=\epsilon(1-\lambda')=5771/10^{13}>2^{-31}.
\]
All parameter inequalities have strict slack. The compact reserved-axis
cost is now $O(V\log p)$, and $\epsilon(1+c)<1$ still holds.

The new scalar mixers use only additions, subtractions and half scalings.
For explicit guard accounting let $M_0=R_0-C$ be the number of additions
in all disjoint-rectangle mixers per invocation. Each pair circuit has
$2(3n-6)$ mixer additions. Let $M$ be their combined mixer count, $V_0$
the number of input-copy incidences, and $J_0$ the output incidences.
The signed schedule uses at most
\[
 T_{\rm scalar}=3v^2(4M+2V_0+4J_0+18v)
\]
elementary scalar operations, including half scalings. The exact checker
verifies $4T_{\rm scalar}+8s+8W+8m<64(W+m+1)^3$ and $2\le s<m^5$.
Thus the retained additive node charge $E=64(W+m+1)^3$ dominates these
operations, per-kernel scalar corrections, and fixed endpoint corrections.
Binary basis changes only
permute coefficient encodings. The generalized stopped-depth proof applies
with the new $W,m,s$, $C_1=5-4\beta+\zeta$ and its explicit $C_0$.
Since $\epsilon C_1<1$, coefficient widths remain $O(p)$.

This conditional construction is integrated in the independent combined
patch listed below, including the new complex
constants, generalized guard and final assembly. The previously published
$2^{-34}$ note and patch remain unchanged.
Neither the computations nor this note provide independent review or a
formal verification of the upstream theorem or the new transfer.
